%!TEX root = ../../csuthesis_main.tex
\chapter{绪论}

\section{研究背景与意义}

\subsection{研究背景}
水下机器人视觉系统在海洋生态监测、水下搜救与海底管线巡检中承担着关键感知任务。然而，水体对光线的选择性吸收与悬浮颗粒的散射会使图像出现严重的色偏、对比度下降与细节模糊。随着水体浑浊度升高，这种退化进一步加剧，直接导致 YOLO 等目标检测模型的漏检率与误检率显著上升。为在安全、可控且可重复的条件下研究这一问题，利用 OpenHUTB 仿真平台调节水体参数并采集多档浑浊度图像，成为一种合理且高效的数据获取途径。

\subsubsection{研究意义}
针对水下图像退化及其对目标识别的影响，本文的研究意义体现在三个层面：在算法层面，系统评估图像增强算法对检测器性能的实际增益，回答“增强是否真的让机器看得更准”这一问题；在数据层面，构建包含清晰、中等、浑浊三档标注图像的水下仿真数据集，为后续对比实验提供基准；在系统层面，将增强与检测流程封装为可交互的 Qt 桌面工具，使算法结果具备可演示性与可复现性。

\section{国内外研究现状}

\subsection{国内研究现状}
国内学者在水下图像增强与目标检测两个方向均开展了大量工作。在增强方面，基于暗通道先验的改进方法被广泛用于水下去雾，研究者通过色彩补偿、背景光估计优化等方式缓解传统暗通道在水下场景中的亮度不足问题\cite{he2009}。在检测方面，多项工作将注意力机制、多尺度特征融合引入 YOLO 框架，以提升水下小目标与遮挡目标的检测精度。但已有研究多将增强与检测作为两个独立环节处理，缺少在统一数据集上对“增强—检测”联合效果的量化对比。

\subsection{国外研究现状}
国外研究更早形成了系统化的综述体系。Yang 等人对水下图像增强与复原方法进行了深入调查，指出物理模型类方法在颜色保真度上具有优势，而数据驱动方法在复杂退化场景中表现出更强的适应性\cite{yang2019}。在增强算法方面，Retinex 理论及其变体被广泛验证，部分学者对基于 Retinex 的水下增强算法进行了嵌入式平台评估，为实时性约束提供了参考\cite{olguin2022}。在检测方面，YOLO 系列从 YOLOv4 到 YOLOv10 均有水下场景的适配研究\cite{yolo2023}。

\subsection{研究述评}
综合来看，现有工作存在两点不足：其一，增强算法的评价多集中于 PSNR、SSIM 等图像质量指标，较少在同一实验设置下检验其对下游检测精度（如 mAP@0.5）的增益；其二，增强模型与检测模型往往分开选型与训练，缺少面向水下机器人实时场景的轻量化协同验证。本文即针对这两点，在 OpenHUTB 仿真数据上开展增强—检测联合对比，并设置单帧 200 ms 的实时性约束。

\section{研究内容与创新点}

\subsection{水下图像数据采集与标注}
基于 OpenHUTB 水下机器人仿真平台采集 800 帧以上图像，按水体浑浊度划分为清晰、中等、浑浊三档，每档不少于 200 帧；完成目标边界框标注不少于 500 例。

\subsection{图像去雾与对比度增强算法研究}
实现并对比暗通道先验、Retinex 及其改进算法，以 PSNR 与 SSIM 为客观指标，要求增强后 PSNR 提升不少于 20 dB、SSIM 不低于 0.80，且单帧处理时间不超过 200 ms。

\subsection{目标识别与对比分析}
采用统一的 YOLO 检测模型，分别在原始图像与增强后图像上训练与测试，对比 mAP@0.5 的变化，量化图像增强对检测性能的提升幅度。

\subsection{系统实现与验证}
基于 Qt 框架实现增强前后图像与检测结果的并排对比显示、指标统计与结果导出，完成系统集成与功能验证。

\section{论文组织架构}

本文共六章。第一章为绪论；第二章介绍水下成像模型、暗通道与 Retinex 原理以及 YOLO 检测基础；第三章给出增强算法流程与实现细节；第四章为增强—检测联合实验与结果分析；第五章为 Qt 系统设计与实现；第六章总结全文并展望后续工作。
